\chapter{Confidence Intervals and Hypothesis Testing}

You have spent the last two chapters learning how to point at a number. Given data,
the method of moments and maximum likelihood hand you an \emph{estimate} --- a single
value $\hat\theta$ that is your best guess for some unknown quantity $\theta$. But a
single number, offered with no hint of how much you should trust it, is nearly
worthless. If a pollster tells you ``$52\%$ of voters favor the measure'' and stops
there, the right response is not ``thank you'' but ``$52\%$ \emph{plus or minus what}?''
A guess without a margin is a guess pretending to be a fact.

This chapter is about attaching a measure of uncertainty to estimates. There are two
great traditions for doing it, and they are really two answers to two different
questions.

\begin{itemize}
  \item \textbf{Confidence intervals} answer: \emph{how wide a net should I cast so that
  the truth is plausibly inside it?} The output is a \emph{range} of values.
  \item \textbf{Hypothesis tests} answer: \emph{is the effect I see real, or could plain
  randomness have produced it?} The output is a \emph{decision} (or a number, the
  p-value, that informs one).
\end{itemize}

Both are built on one idea you already own from Chapter~21: the \emph{sampling
distribution}. Everything in this chapter --- every $\pm$, every p-value, every t-test
--- is an accountant's report on how much an estimate wobbles when you redraw the
data. We will build that idea carefully, then watch it grow into the two most-used (and
two most-\emph{misused}) tools in all of applied statistics, paying close attention to
what they do and, just as importantly, what they do \emph{not} do.

\section{The sampling distribution: the engine underneath}
\label{ch32:sampling}

Imagine the unknown quantity is the mean height $\mu$ of adult women in a country. You
cannot measure everyone, so you draw a random sample of $n$ women and compute the
sample mean
\[
  \bar X_n = \frac{1}{n}\sum_{i=1}^n X_i .
\]
Here is the conceptual leap that makes all of inference work. Your particular sample
gave you, say, $\bar X_n = 164.3$ cm. But that number was an accident of \emph{which}
$n$ women happened to be chosen. A different random sample would have given a slightly
different mean. So $\bar X_n$ is not a fixed number --- it is itself a random variable,
with its own distribution. That distribution is called the \textbf{sampling
distribution} of the estimator.

\begin{intuitionbox}
Picture a giant, hypothetical experiment: draw a sample of size $n$, write down
$\bar X_n$, throw the sample back, and repeat forever. The histogram of all those
$\bar X_n$ values \emph{is} the sampling distribution. We never actually run this
experiment --- we have only the one sample we paid for. But probability theory lets us
\emph{predict} the shape of that histogram without running it. That prediction is the
whole game.
\end{intuitionbox}

For the sample mean we know the prediction exactly. If the $X_i$ are i.i.d.\ with mean
$\mu$ and variance $\sigma^2$, then by linearity and independence (Chapters 8--9),
\[
  \E[\bar X_n] = \mu, \qquad \Var(\bar X_n) = \frac{\sigma^2}{n}.
\]
The estimator is centered on the truth (it is \emph{unbiased}) and its variance shrinks
like $1/n$. The standard deviation of an estimator has a special name: the
\textbf{standard error}.
\[
  \SD(\bar X_n) = \frac{\sigma}{\sqrt n} \;=:\; \mathrm{SE}.
\]

\begin{trapbox}
Do not confuse the \emph{standard deviation} $\sigma$ (how spread out individual people
are) with the \emph{standard error} $\sigma/\sqrt n$ (how spread out the \emph{sample
mean} is across hypothetical repetitions). Heights vary by maybe $7$ cm person to
person; the mean of $1000$ heights varies by only $7/\sqrt{1000}\approx 0.22$ cm. As
$n$ grows, people stay just as different but the \emph{average} becomes razor-sharp.
This is the single most common confusion in introductory statistics.
\end{trapbox}

And the shape? This is where the Central Limit Theorem (Chapter~21) does its magic. No
matter what the distribution of an individual $X_i$ looks like --- skewed, bimodal,
discrete --- the sampling distribution of $\bar X_n$ is approximately Normal for large
$n$:
\[
  \bar X_n \;\overset{\text{approx}}{\sim}\; \Normal\!\left(\mu,\ \frac{\sigma^2}{n}\right),
  \qquad\text{equivalently}\qquad
  \frac{\bar X_n - \mu}{\sigma/\sqrt n} \;\dto\; \Normal(0,1).
\]

\begin{whybox}
Why is the Normal curve the universal shape of sampling distributions? Because
$\bar X_n$ is an \emph{average of many small independent contributions}, and the CLT
says such averages forget the idiosyncrasies of their ingredients and converge to the
same bell curve. This is the deep reason confidence intervals and z-tests work for
almost any underlying population: we are never really using the population's shape, only
the shape of the \emph{average}, and that shape is always Normal in the limit.
\end{whybox}

Figure~\ref{ch32:fig-sampling} shows the punchline visually: a wildly non-Normal
population (a skewed exponential) whose sample means nonetheless pile up into a tidy
bell as $n$ grows, narrowing all the while.

\begin{figure}[H]
\centering
\begin{tikzpicture}
\begin{axis}[
  width=12cm, height=6.2cm,
  xlabel={value of $\bar X_n$}, ylabel={density},
  domain=0:3, samples=160,
  legend pos=north east, legend style={font=\small},
  ymin=0, axis lines=left, clip=false,
]
% population mean = 1; SE = 1/sqrt(n); CLT bells of increasing n
\addplot[densityred, thick] {1/(1*0.7071)*0.39894*exp(-((x-1)^2)/(2*0.5))};
\addlegendentry{$n=2$}
\addplot[quotientorange, thick] {1/(1*0.3162)*0.39894*exp(-((x-1)^2)/(2*0.1))};
\addlegendentry{$n=10$}
\addplot[distblue, thick] {1/(1*0.1414)*0.39894*exp(-((x-1)^2)/(2*0.02))};
\addlegendentry{$n=50$}
\draw[gray,dashed] (axis cs:1,0) -- (axis cs:1,3.1);
\node[gray,anchor=south] at (axis cs:1,3.1) {\small truth $\mu=1$};
\end{axis}
\end{tikzpicture}
\caption{Sampling distributions of $\bar X_n$ for samples drawn from a skewed
population with mean $\mu=1$. Two things happen as $n$ grows: the curves become more
Normal (CLT) and narrower (standard error $\sigma/\sqrt n \to 0$). All are centered on
the truth.}
\label{ch32:fig-sampling}
\end{figure}

\begin{happeningbox}
Inference is a kind of reverse engineering. We observe one point drawn from the
sampling distribution (our single $\bar X_n$) and we ask: which values of the unknown
$\mu$ would make this observation unsurprising? A confidence interval is the set of
such $\mu$'s. A hypothesis test fixes one candidate $\mu$ and asks whether our
observation is surprising under it. Same engine, two readouts.
\end{happeningbox}

\section{Confidence intervals}
\label{ch32:ci}

\subsection{Building one from the sampling distribution}

Start with the clean case: i.i.d.\ data, large $n$, and suppose for a moment we know
$\sigma$. Standardize the sample mean:
\[
  Z = \frac{\bar X_n - \mu}{\sigma/\sqrt n} \approx \Normal(0,1).
\]
From the standard Normal we know that $95\%$ of the probability lies between $-1.96$ and
$+1.96$. Write $z_{\alpha/2}$ for the value with $\alpha/2$ probability in the upper
tail; for $95\%$ confidence, $\alpha = 0.05$ and $z_{0.025} = 1.96$. Then
\[
  \P\!\left(-z_{\alpha/2} \le \frac{\bar X_n - \mu}{\sigma/\sqrt n} \le z_{\alpha/2}\right)
  \approx 1-\alpha .
\]
Now do nothing but algebra --- multiply through by the standard error and rearrange to
isolate $\mu$ in the middle:
\[
  \P\!\left(\bar X_n - z_{\alpha/2}\frac{\sigma}{\sqrt n}
  \ \le\ \mu\ \le\
  \bar X_n + z_{\alpha/2}\frac{\sigma}{\sqrt n}\right) \approx 1-\alpha .
\]

\begin{formalbox}
A \textbf{$(1-\alpha)$ confidence interval} for $\mu$ (known $\sigma$, large $n$) is
\[
  \boxed{\ \bar X_n \ \pm\ z_{\alpha/2}\,\frac{\sigma}{\sqrt n}\ } .
\]
The half-width $z_{\alpha/2}\,\sigma/\sqrt n$ is the \textbf{margin of error}. It shrinks
with $\sqrt n$: to halve the margin you must \emph{quadruple} the sample.
\end{formalbox}

\begin{example}[A polling margin]
\label{ch32:ex-poll}
A poll of $n=1000$ voters finds $\hat p = 0.52$ in favor. For a proportion the variance
of a single response is $p(1-p)$, so the standard error is $\sqrt{p(1-p)/n}$; plugging
in $\hat p$ gives $\sqrt{0.52\cdot 0.48/1000}\approx 0.0158$. The $95\%$ interval is
\[
  0.52 \pm 1.96 \times 0.0158 = 0.52 \pm 0.031 = (0.489,\ 0.551).
\]
This is the famous ``margin of error of about $3$ points.'' Note the interval includes
$0.50$: on this evidence we cannot rule out a tie. The headline ``leading $52$ to $48$''
oversells what the data support.
\end{example}

\subsection{What a confidence interval really means --- and does not}

This is the most misunderstood concept in statistics, so let us be slow and exact.

\begin{trapbox}
\textbf{The interval $(0.489, 0.551)$ does NOT mean ``there is a $95\%$ probability that
$\mu$ lies in $(0.489, 0.551)$.''} In the frequentist framework $\mu$ is a fixed,
unknown constant --- not random. It is either in that particular interval or it is not;
the probability is $0$ or $1$, we just do not know which. The randomness was in the
\emph{data}, which produced the \emph{endpoints}, not in $\mu$.
\end{trapbox}

So what is the $95\%$ attached to? It is a property of the \emph{procedure}, not of any
single interval.

\begin{focusbox}
\textbf{The correct interpretation.} The recipe ``compute $\bar X_n \pm 1.96\,
\mathrm{SE}$'' is a machine that spits out a random interval every time you feed it a
fresh sample. \emph{Across many independent samples, $95\%$ of the intervals this machine
produces will contain the true $\mu$.} Your one interval is a single draw from that
machine. You do not know if it is one of the lucky $95\%$ or the unlucky $5\%$ --- but
you trust the machine because it is right $95\%$ of the time.
\end{focusbox}

Figure~\ref{ch32:fig-coverage} makes this tangible. We simulate twenty samples, draw the
twenty intervals, and mark the fixed truth with a vertical line. Most intervals cover
it; a stray one or two miss. The ``$95\%$'' is the long-run \emph{coverage rate}, the
fraction of horizontal bars that cross the line.

\begin{figure}[H]
\centering
\begin{tikzpicture}
\begin{axis}[
  width=12cm, height=7cm,
  xlabel={estimate scale}, ylabel={sample number},
  xmin=-3.2, xmax=4.2, ymin=0, ymax=21,
  axis lines=left, ytick=\empty,
]
\draw[densityred, thick, dashed] (axis cs:0,0) -- (axis cs:0,21);
\node[densityred, anchor=south] at (axis cs:0,21) {\small true $\mu$};
% each interval: center +- 1.5 (the "95% half-width" on this scale); 19 cover, 1 misses.
% covering intervals (blue), each centered at \c with half-width 1.5:
\draw[distblue, very thick] (axis cs:-1.9,1) -- (axis cs:1.1,1); \fill[distblue] (axis cs:-0.4,1) circle (1.3pt);
\draw[distblue, very thick] (axis cs:-0.9,2) -- (axis cs:2.1,2); \fill[distblue] (axis cs:0.6,2) circle (1.3pt);
\draw[distblue, very thick] (axis cs:-2.4,3) -- (axis cs:0.6,3); \fill[distblue] (axis cs:-0.9,3) circle (1.3pt);
\draw[distblue, very thick] (axis cs:-1.2,4) -- (axis cs:1.8,4); \fill[distblue] (axis cs:0.3,4) circle (1.3pt);
\draw[distblue, very thick] (axis cs:-0.5,5) -- (axis cs:2.5,5); \fill[distblue] (axis cs:1.0,5) circle (1.3pt);
\draw[distblue, very thick] (axis cs:-2.0,6) -- (axis cs:1.0,6); \fill[distblue] (axis cs:-0.5,6) circle (1.3pt);
\draw[distblue, very thick] (axis cs:-0.7,7) -- (axis cs:2.3,7); \fill[distblue] (axis cs:0.8,7) circle (1.3pt);
\draw[distblue, very thick] (axis cs:-2.6,8) -- (axis cs:0.4,8); \fill[distblue] (axis cs:-1.1,8) circle (1.3pt);
\draw[distblue, very thick] (axis cs:-1.3,9) -- (axis cs:1.7,9); \fill[distblue] (axis cs:0.2,9) circle (1.3pt);
\draw[distblue, very thick] (axis cs:-1.0,10) -- (axis cs:2.0,10); \fill[distblue] (axis cs:0.5,10) circle (1.3pt);
\draw[distblue, very thick] (axis cs:-2.2,11) -- (axis cs:0.8,11); \fill[distblue] (axis cs:-0.7,11) circle (1.3pt);
\draw[distblue, very thick] (axis cs:-0.6,12) -- (axis cs:2.4,12); \fill[distblue] (axis cs:0.9,12) circle (1.3pt);
\draw[distblue, very thick] (axis cs:-1.8,13) -- (axis cs:1.2,13); \fill[distblue] (axis cs:-0.3,13) circle (1.3pt);
\draw[distblue, very thick] (axis cs:-1.1,14) -- (axis cs:1.9,14); \fill[distblue] (axis cs:0.4,14) circle (1.3pt);
\draw[distblue, very thick] (axis cs:-2.1,15) -- (axis cs:0.9,15); \fill[distblue] (axis cs:-0.6,15) circle (1.3pt);
\draw[distblue, very thick] (axis cs:-0.8,16) -- (axis cs:2.2,16); \fill[distblue] (axis cs:0.7,16) circle (1.3pt);
\draw[distblue, very thick] (axis cs:-1.7,17) -- (axis cs:1.3,17); \fill[distblue] (axis cs:-0.2,17) circle (1.3pt);
\draw[distblue, very thick] (axis cs:-0.4,18) -- (axis cs:2.6,18); \fill[distblue] (axis cs:1.1,18) circle (1.3pt);
\draw[distblue, very thick] (axis cs:-2.3,19) -- (axis cs:0.7,19); \fill[distblue] (axis cs:-0.8,19) circle (1.3pt);
% the one missing interval (red): lies entirely to the right of the truth
\draw[trapred, very thick] (axis cs:0.9,20) -- (axis cs:3.9,20); \fill[trapred] (axis cs:2.4,20) circle (1.3pt);
\end{axis}
\end{tikzpicture}
\caption{Twenty $95\%$ confidence intervals from twenty independent samples. The truth
$\mu$ (red dashed line) is fixed; the intervals jiggle. Nineteen of the twenty cover the
truth (blue); one misses (red). Over the long run, $95\%$ cover. The confidence level
describes the \emph{procedure}, not your particular interval.}
\label{ch32:fig-coverage}
\end{figure}

\begin{historybox}
The confidence interval was introduced by Jerzy Neyman in a landmark 1937 paper, ``Outline
of a Theory of Statistical Estimation Based on the Classical Theory of Probability.''
Neyman was deliberately, almost stubbornly, careful: he insisted the probability
statement attaches to the \emph{random interval before the data are seen}, never to the
fixed parameter after. He knew the temptation to say ``$95\%$ chance $\mu$ is in here''
and built the whole theory to resist it. The misinterpretation has persisted ever since,
a sign of how badly people \emph{want} the Bayesian statement --- which, as
Chapter~31 shows, is exactly what a Bayesian credible interval delivers, at the price of
a prior.
\end{historybox}

\begin{connectionbox}
The frequentist confidence interval and the Bayesian credible interval (Chapter~31)
often produce nearly identical numbers but answer different questions. The credible
interval \emph{does} let you say ``$95\%$ probability $\mu$ is in here,'' because in the
Bayesian world $\mu$ has a probability distribution. The confidence interval refuses
that statement and instead guarantees long-run coverage. With a flat prior and a Normal
likelihood the two intervals coincide algebraically --- a reassuring meeting point of
the two great schools.
\end{connectionbox}

\subsection{When you do not know $\sigma$: the t-interval}

We assumed we knew the population standard deviation $\sigma$. In real life we never do
--- we estimate it from the same data with the \emph{sample standard deviation}
\[
  s = \sqrt{\frac{1}{n-1}\sum_{i=1}^n (X_i - \bar X_n)^2}.
\]
The $n-1$ (rather than $n$) makes $s^2$ an unbiased estimator of $\sigma^2$; we lose one
degree of freedom because we already spent the data to estimate $\bar X_n$. Substituting
$s$ for $\sigma$ introduces extra uncertainty --- now \emph{both} the numerator and the
denominator of our standardized statistic are random. William Gosset worked out the
consequence: the statistic
\[
  T = \frac{\bar X_n - \mu}{s/\sqrt n}
\]
no longer follows a Normal distribution but \textbf{Student's $t$-distribution} with
$n-1$ degrees of freedom. The $t$ looks like the Normal but with heavier tails ---
fatter, because the extra randomness in $s$ makes large deviations more likely.

\begin{figure}[H]
\centering
\begin{tikzpicture}
\begin{axis}[
  width=12cm, height=5.6cm,
  xlabel={$t$}, ylabel={density}, domain=-4:4, samples=160,
  legend pos=north east, legend style={font=\small}, ymin=0, axis lines=left,
]
\addplot[distblue, thick] {0.39894*exp(-x^2/2)};
\addlegendentry{Normal}
% t with 3 df: c=0.36755; (1+x^2/3)^{-2}
\addplot[densityred, thick] {0.36755*(1+x^2/3)^(-2)};
\addlegendentry{$t_3$ (heavier tails)}
\end{axis}
\end{tikzpicture}
\caption{Student's $t$ with $3$ degrees of freedom versus the standard Normal. The $t$
has heavier tails: estimating $\sigma$ from the data costs you, so you must reach
further out for the same confidence. As $n\to\infty$, $s\to\sigma$ and $t\to$ Normal.}
\label{ch32:fig-t}
\end{figure}

The confidence interval is identical in form, with the critical value taken from the
$t$-table instead of the Normal:
\[
  \boxed{\ \bar X_n \ \pm\ t_{n-1,\,\alpha/2}\,\frac{s}{\sqrt n}\ } .
\]
Because $t_{n-1,\alpha/2} > z_{\alpha/2}$, the $t$-interval is always a little
\emph{wider} --- the price of not knowing $\sigma$. For $n = 10$, the $95\%$
critical value is $t_{9,0.025} = 2.262$, noticeably more than $1.96$. By $n = 100$ it
has fallen to $1.984$, almost the Normal value. For large samples the distinction
vanishes, which is why the z- and t-intervals merge.

\begin{historybox}
Student's $t$ is named for a pseudonym. William Sealy Gosset derived it in 1908 while
working as a chemist at the Guinness brewery in Dublin, where he needed to draw
conclusions from \emph{small} batches of barley and yeast --- samples far too small for
the large-$n$ Normal approximation. Guinness forbade employees from publishing (a rival
had leaked trade secrets), so Gosset published under the modest pen name ``Student.''
The most-used distribution in applied statistics is thus named after a man hiding from
his employer. R.~A.~Fisher later put the result on rigorous footing and gave it the
letter $t$.
\end{historybox}

\begin{example}[A small-sample t-interval]
\label{ch32:ex-tinterval}
Six measurements of a chemical yield: $48, 51, 49, 53, 50, 52$ (grams). The sample mean
is $\bar X = 50.5$. The deviations are $-2.5, 0.5, -1.5, 2.5, -0.5, 1.5$, whose squares
sum to $6.25+0.25+2.25+6.25+0.25+2.25 = 17.5$, so
\[
  s^2 = \frac{17.5}{6-1} = 3.5,\qquad s = 1.871,\qquad \frac{s}{\sqrt 6} = 0.764.
\]
With $5$ degrees of freedom, $t_{5,0.025} = 2.571$. The $95\%$ interval is
\[
  50.5 \pm 2.571 \times 0.764 = 50.5 \pm 1.96 = (48.54,\ 52.46)\ \text{grams}.
\]
Had we naively used $1.96$ instead of $2.571$, we would have reported a falsely narrow
interval --- overstating our precision, which is exactly the error the $t$ corrects.
\end{example}

\section{Hypothesis testing}
\label{ch32:testing}

A confidence interval casts a net. A hypothesis test asks a sharper, binary question:
\emph{is there any effect at all, or could what I see be pure chance?} The logic is a
courtroom's --- presumption of innocence, evidence, a verdict ``beyond reasonable
doubt'' --- dressed in the language of probability.

\subsection{The null, the alternative, and the logic of contradiction}

We set up two competing claims:

\begin{itemize}
  \item The \textbf{null hypothesis} $H_0$: the boring, status-quo, ``nothing is going
  on'' claim. The drug has no effect; the coin is fair; the two web pages convert at the
  same rate. $H_0$ is the defendant, presumed innocent.
  \item The \textbf{alternative hypothesis} $H_1$ (or $H_a$): the interesting claim we
  would believe only if the evidence forces us. The drug works; the coin is biased; page
  B is better.
\end{itemize}

The strategy is \emph{proof by contradiction, softened by probability}. We provisionally
assume $H_0$ is true. We then ask: \emph{if the null were true, how surprising would my
data be?} If the data are wildly unlikely under $H_0$, we take that as evidence against
$H_0$ and \emph{reject} it. If the data are unremarkable under $H_0$, we have no grounds
to abandon it and we \emph{fail to reject}.

\begin{trapbox}
We say ``fail to reject $H_0$,'' never ``accept $H_0$.'' Absence of evidence is not
evidence of absence. A jury that does not convict declares ``not guilty,'' not
``innocent.'' Failing to find an effect could mean there is none --- or that your sample
was too small to see it. The asymmetry is deliberate and important.
\end{trapbox}

\subsection{Test statistics and the null distribution}

To measure ``how surprising,'' we compress the data into a single \textbf{test
statistic} whose distribution we know \emph{assuming $H_0$ is true}. For testing a mean
with known $\sigma$, the natural statistic is again the standardized sample mean,
\[
  Z = \frac{\bar X_n - \mu_0}{\sigma/\sqrt n},
\]
where $\mu_0$ is the value claimed by $H_0$. Under $H_0$ this $Z$ is approximately
$\Normal(0,1)$ --- this is the \textbf{null distribution}. A $Z$ near $0$ is what the
null predicts; a $Z$ far out in the tails is the kind of thing the null almost never
produces, and so it indicts the null.

\subsection{The p-value}

\begin{formalbox}
The \textbf{p-value} is the probability, \emph{computed assuming $H_0$ is true}, of
observing a test statistic at least as extreme as the one you actually got:
\[
  p = \P\big(\,\text{statistic at least as extreme as observed}\ \big|\ H_0\,\big).
\]
Small $p$ means your data would be a rare fluke if the null held --- evidence
\emph{against} $H_0$.
\end{formalbox}

For a two-sided test with observed value $z_{\text{obs}}$, ``as extreme'' means as far
from $0$ in either direction, so
\[
  p = \P(|Z| \ge |z_{\text{obs}}|) = 2\big(1 - \Phi(|z_{\text{obs}}|)\big),
\]
where $\Phi$ is the standard Normal CDF. Figure~\ref{ch32:fig-pvalue} shows the p-value
as the shaded tail area.

\begin{figure}[H]
\centering
\begin{tikzpicture}
\begin{axis}[
  width=12cm, height=5.6cm,
  xlabel={$z$ (test statistic under $H_0$)}, ylabel={density},
  domain=-4:4, samples=200, ymin=0, axis lines=left, clip=false,
]
\addplot[distblue, thick] {0.39894*exp(-x^2/2)};
% right tail beyond 2.1
\addplot[draw=none, fill=trapred!45, domain=2.1:4] {0.39894*exp(-x^2/2)} \closedcycle;
% left tail beyond -2.1
\addplot[draw=none, fill=trapred!45, domain=-4:-2.1] {0.39894*exp(-x^2/2)} \closedcycle;
\draw[densityred, thick] (axis cs:2.1,0) -- (axis cs:2.1,0.34);
\node[densityred, anchor=south, font=\small] at (axis cs:2.1,0.34) {$z_{\text{obs}}$};
\node[trapred, anchor=west, font=\small] at (axis cs:2.6,0.07) {tail};
\node[trapred, anchor=east, font=\small] at (axis cs:-2.6,0.07) {tail};
\node[anchor=south, font=\small] at (axis cs:0,0.40) {$p = $ both shaded tails};
\end{axis}
\end{tikzpicture}
\caption{The p-value for a two-sided z-test is the total tail area beyond the observed
statistic (here $z_{\text{obs}}=2.1$), in \emph{both} directions. It is the probability
the null would generate something at least this extreme by chance alone.}
\label{ch32:fig-pvalue}
\end{figure}

We compare $p$ to a pre-chosen threshold $\alpha$ (the \textbf{significance level},
conventionally $0.05$). If $p \le \alpha$, we reject $H_0$ and call the result
``statistically significant.'' If $p > \alpha$, we fail to reject.

\begin{example}[A z-test for a fair coin]
\label{ch32:ex-coin}
You flip a coin $n=100$ times and see $62$ heads. Is it fair? Let $H_0\!: p = 0.5$ versus
$H_1\!: p \ne 0.5$. Under $H_0$ the number of heads has mean $50$ and standard deviation
$\sqrt{100 \cdot 0.5 \cdot 0.5}=5$. The standardized statistic is
\[
  z_{\text{obs}} = \frac{62 - 50}{5} = 2.4.
\]
The two-sided p-value is $p = 2(1-\Phi(2.4)) = 2(0.0082) = 0.0164$. Since $0.0164 <
0.05$, we reject $H_0$: this coin looks biased. Had we seen $58$ heads, $z = 1.6$,
$p = 0.11 > 0.05$, and we would not reject --- $58$ heads is the sort of thing a fair
coin throws up often enough.
\end{example}

\subsection{What the p-value is not}

The p-value may be the single most abused number in science. Three corrections:

\begin{trapbox}
\textbf{The p-value is NOT the probability that $H_0$ is true.} It is computed
\emph{assuming} $H_0$ is true: $p = \P(\text{data}\given H_0)$, not $\P(H_0\given
\text{data})$. Confusing these two is the \emph{prosecutor's fallacy} --- swapping the
direction of a conditional probability (Chapter~5). A small p-value says ``the data are
surprising if $H_0$ holds,'' which is genuine evidence against $H_0$, but it is not a
posterior probability. Getting from $\P(\text{data}\given H_0)$ to $\P(H_0\given
\text{data})$ requires Bayes' theorem and a prior --- which the p-value does not have.
\end{trapbox}

\begin{trapbox}
\textbf{Statistical significance is not practical importance.} With a large enough $n$,
\emph{any} effect, however trivially small, will eventually produce $p<0.05$ --- because
the standard error $\sigma/\sqrt n$ shrinks to zero and even a microscopic difference
lands many standard errors from the null. A drug that lowers blood pressure by
$0.1$ mmHg can be ``highly significant'' in a study of a million people and utterly
useless in a clinic. Always report the \emph{effect size} (and a confidence interval),
not just the p-value.
\end{trapbox}

\begin{trapbox}
\textbf{$p = 0.049$ and $p = 0.051$ are essentially the same evidence.} The $0.05$
cutoff is a convention, not a law of nature (Fisher proposed it almost offhandedly).
Treating $0.05$ as a bright line between ``real'' and ``not real'' --- and worse,
running analyses until you crawl under it (``p-hacking'') --- has fueled a replication
crisis across the sciences. The p-value is a continuous measure of surprise; report it
as such.
\end{trapbox}

\begin{paradoxbox}
\textbf{The base-rate trap behind ``significant'' findings.} Suppose only $1$ in $10$
hypotheses a lab tests is actually true, and tests run at $\alpha=0.05$ with power
$0.8$ (defined below). Out of $1000$ tests, $100$ effects are real and $900$ are null.
The tests flag $0.8\times 100 = 80$ true effects and $0.05\times 900 = 45$ false ones.
So among the $125$ ``significant'' results, $45/125 = 36\%$ are \emph{false positives}
--- despite every individual test having only a $5\%$ error rate! This is Bayes' theorem
again (Chapter~5): a low false-positive \emph{rate} per test does not imply a low
false-discovery \emph{proportion} when true effects are rare. It explains why so many
published ``significant'' findings fail to replicate.
\end{paradoxbox}

\section{Two kinds of error, and the power to detect}
\label{ch32:errors}

Any decision rule that says ``reject'' or ``do not reject'' can be wrong in two ways.
The truth ($H_0$ or $H_1$) crossed with our verdict gives a $2\times 2$ table:

\begin{center}
\renewcommand{\arraystretch}{1.4}
\begin{tabular}{l|cc}
 & \textbf{$H_0$ true} & \textbf{$H_1$ true} \\ \hline
\textbf{Reject $H_0$} & Type I error ($\alpha$) & correct (power) \\
\textbf{Fail to reject} & correct & Type II error ($\beta$) \\
\end{tabular}
\end{center}

\begin{itemize}
  \item A \textbf{Type I error} is convicting the innocent: rejecting $H_0$ when it is
  true. We \emph{control} its probability directly --- it is exactly $\alpha$, the
  significance level we chose. Setting $\alpha=0.05$ means ``I am willing to cry wolf
  $5\%$ of the time when there is no wolf.''
  \item A \textbf{Type II error} is acquitting the guilty: failing to reject $H_0$ when
  $H_1$ is true. Its probability is called $\beta$.
  \item The \textbf{power} of a test is $1-\beta$: the probability of correctly detecting
  a real effect. More power is better.
\end{itemize}

\begin{intuitionbox}
Think of $\alpha$ as the sensitivity of a smoke alarm. Crank it up (large $\alpha$,
quick to reject) and you catch every fire (high power) but also scream at burnt toast
(many Type I errors). Quiet it down (small $\alpha$) and you stop the false alarms but
risk sleeping through a real blaze (Type II errors). For \emph{fixed} $n$ there is a
genuine trade-off: pushing $\alpha$ down pushes $\beta$ up. The only way to beat the
trade-off --- to shrink both errors at once --- is to gather more data.
\end{intuitionbox}

Figure~\ref{ch32:fig-power} shows the two error regions as overlapping bells: the null
distribution centered at $\mu_0$ and the (true) alternative centered at some $\mu_1$.
The cutoff $c$ slices both. The right tail of the null beyond $c$ is $\alpha$; the left
tail of the alternative below $c$ is $\beta$.

\begin{figure}[H]
\centering
\begin{tikzpicture}
\begin{axis}[
  width=13cm, height=6cm,
  xlabel={test statistic / sample mean}, ylabel={density},
  domain=-4:7, samples=200, ymin=0, axis lines=left, clip=false,
  xtick={0,3}, xticklabels={$\mu_0$,$\mu_1$},
]
% null centered 0, alt centered 3, cutoff at 1.9
\addplot[distblue, thick] {0.39894*exp(-x^2/2)};
\addplot[targetgreen, thick] {0.39894*exp(-(x-3)^2/2)};
% alpha: right tail of null beyond 1.9
\addplot[draw=none, fill=trapred!50, domain=1.9:4] {0.39894*exp(-x^2/2)} \closedcycle;
% beta: left tail of alt below 1.9
\addplot[draw=none, fill=quotientorange!45, domain=-1:1.9] {0.39894*exp(-(x-3)^2/2)} \closedcycle;
\draw[densityred, thick] (axis cs:1.9,0) -- (axis cs:1.9,0.42);
\node[densityred, anchor=south, font=\small] at (axis cs:1.9,0.42) {cutoff $c$};
\node[distblue, anchor=south, font=\small] at (axis cs:0,0.40) {$H_0$};
\node[targetgreen, anchor=south, font=\small] at (axis cs:3,0.40) {$H_1$};
\node[trapred, font=\small, anchor=west] at (axis cs:2.3,0.045) {$\alpha$};
\node[quotientorange, font=\small, anchor=east] at (axis cs:1.6,0.05) {$\beta$};
\end{axis}
\end{tikzpicture}
\caption{Type I and Type II errors. Reject when the statistic exceeds the cutoff $c$.
The red area $\alpha$ is the chance of a false alarm under $H_0$; the orange area
$\beta$ is the chance of missing a real effect under $H_1$. Power is the area under the green curve to
the \emph{right} of $c$, i.e.\ $1-\beta$. Moving $c$ left trades $\beta$ for $\alpha$.}
\label{ch32:fig-power}
\end{figure}

\begin{whybox}
Why does power rise with sample size? Larger $n$ shrinks the standard error
$\sigma/\sqrt n$, so \emph{both} bells in Figure~\ref{ch32:fig-power} narrow. Two narrow
bells centered at $\mu_0$ and $\mu_1$ overlap far less than two wide ones. Less overlap
means the cutoff can keep $\alpha$ fixed while slicing off much less of the alternative
--- so $\beta$ falls and power climbs. This is the mathematical reason ``get more data''
is the universal cure: it separates the hypotheses.
\end{whybox}

\begin{example}[Power and sample-size planning]
\label{ch32:ex-power}
You want to detect whether a new process raises mean yield from $\mu_0=50$ to
$\mu_1=52$, with $\sigma=4$, at $\alpha=0.05$ (one-sided). The cutoff in raw units is
$c = \mu_0 + 1.645\,\sigma/\sqrt n$. To achieve power $0.8$ you need $\mu_1$ to sit
$0.84$ standard errors above $c$ (since $\Phi(0.84)=0.80$). Setting
$\mu_1 = c + 0.84\,\sigma/\sqrt n$ and solving,
\[
  \mu_1 - \mu_0 = (1.645 + 0.84)\frac{\sigma}{\sqrt n}
  \ \Rightarrow\
  n = \left(\frac{(1.645+0.84)\,\sigma}{\mu_1-\mu_0}\right)^2
  = \left(\frac{2.485 \times 4}{2}\right)^2 \approx 24.7.
\]
So $n=25$ samples give roughly $80\%$ power. Halve the effect you wish to detect and $n$
\emph{quadruples} --- the same $\sqrt n$ law that governs confidence intervals. Power
analysis done \emph{before} collecting data is how careful experimenters avoid the
heartbreak of an underpowered study that could never have found the effect.
\end{example}

\section{The duality of intervals and tests}
\label{ch32:duality}

Confidence intervals and hypothesis tests are not two separate tools --- they are two
views of the same picture, and seeing the link makes both clearer.

\begin{focusbox}
\textbf{Interval--test duality.} A two-sided level-$\alpha$ test of $H_0\!:\mu=\mu_0$
rejects exactly when $\mu_0$ falls \emph{outside} the $(1-\alpha)$ confidence interval.
Equivalently: \emph{the confidence interval is the set of all null values $\mu_0$ that
would not be rejected.}
\end{focusbox}

The algebra is immediate. The test fails to reject when
$\left|\frac{\bar X_n - \mu_0}{\sigma/\sqrt n}\right| \le z_{\alpha/2}$, which rearranges
to $\bar X_n - z_{\alpha/2}\,\sigma/\sqrt n \le \mu_0 \le \bar X_n + z_{\alpha/2}\,
\sigma/\sqrt n$ --- precisely the statement ``$\mu_0$ is inside the interval.''
Figure~\ref{ch32:fig-duality} draws the equivalence: a null value lands outside the bar
exactly when the test rejects it.

\begin{figure}[H]
\centering
\begin{tikzpicture}[x=1.4cm,y=1cm]
% number line
\draw[->,gray] (-3.4,0) -- (3.6,0) node[right,black,font=\small] {$\mu_0$};
% the confidence interval as a thick bar
\draw[distblue, line width=3pt] (-1.96,1.1) -- (1.96,1.1);
\fill[distblue] (0,1.1) circle (2pt);
% endpoint ticks down to the axis
\draw[distblue,dashed] (-1.96,0) -- (-1.96,1.1);
\draw[distblue,dashed] (1.96,0) -- (1.96,1.1);
\node[below,font=\footnotesize] at (-1.96,0) {$\bar X_n - z_{\alpha/2}\mathrm{SE}$};
\node[below,font=\footnotesize] at (0,0) {$\bar X_n$};
\node[below,font=\footnotesize] at (1.96,0) {$\bar X_n + z_{\alpha/2}\mathrm{SE}$};
\node[distblue, anchor=south, font=\small] at (0,1.25) {$(1-\alpha)$ confidence interval};
\node[targetgreen, font=\footnotesize, anchor=north] at (0,0.95) {test fails to reject these $\mu_0$};
% a rejected null on the right
\fill[trapred] (2.8,1.1) circle (2.2pt);
\node[trapred, font=\footnotesize, anchor=south] at (2.8,1.25) {$\mu_0$ rejected};
\draw[trapred,<-] (2.05,1.1) -- (2.7,1.1);
\end{tikzpicture}
\caption{Interval--test duality. The confidence interval (blue bar) is exactly the set
of null values a two-sided level-$\alpha$ test would \emph{not} reject. A candidate
$\mu_0$ inside the bar survives the test; one outside it (red) is rejected. One picture,
two tools.}
\label{ch32:fig-duality}
\end{figure}

This duality pays off in practice. Reporting a confidence interval is strictly more
informative than reporting a single test: it tells you \emph{every} null you would
reject at once, plus the plausible range and the effect size --- which is exactly why
statisticians increasingly prefer intervals to bare p-values.

\section{The t-test in practice}
\label{ch32:ttest}

Combining the pieces: when $\sigma$ is unknown (always, in practice) and we test a mean,
we use the \textbf{t-test}. The statistic is
\[
  T = \frac{\bar X_n - \mu_0}{s/\sqrt n},
\]
compared against the $t_{n-1}$ distribution. There are three common flavors:

\begin{itemize}
  \item \textbf{One-sample t-test}: is the mean of one group equal to a fixed $\mu_0$?
  \item \textbf{Two-sample t-test}: do two independent groups have equal means? The
  statistic is $T = (\bar X - \bar Y)/\mathrm{SE}_{\text{diff}}$, where the standard
  error of the difference combines both groups' variability.
  \item \textbf{Paired t-test}: when observations come in matched pairs (before/after on
  the same subject), test whether the mean \emph{difference} is zero --- a one-sample
  test on the differences, which cancels person-to-person variation.
\end{itemize}

\begin{example}[Two-sample t-test]
\label{ch32:ex-twosample}
A field gets $\bar X = 52$ bushels/acre with fertilizer ($n_X=10$, $s_X=4$); a control
gets $\bar Y = 48$ ($n_Y=10$, $s_Y=5$). Does fertilizer help? With $H_0\!:\mu_X=\mu_Y$,
the standard error of the difference (assuming independent groups) is
\[
  \mathrm{SE}_{\text{diff}} = \sqrt{\frac{s_X^2}{n_X}+\frac{s_Y^2}{n_Y}}
  = \sqrt{\frac{16}{10}+\frac{25}{10}} = \sqrt{4.1} = 2.025.
\]
The statistic is $T = (52-48)/2.025 = 1.98$. With roughly $17$ degrees of freedom (a
standard approximation for unequal variances), the two-sided critical value is about
$2.11$. Since $1.98 < 2.11$, $p \approx 0.064$: we (just barely) fail to reject at
$\alpha=0.05$. The fertilizer \emph{looks} helpful, but with only ten plots per group we
cannot rule out chance. The accurate report is: ``estimated gain $4$ bushels, $95\%$ CI
roughly $(-0.3, 8.3)$, $p=0.06$'' --- not ``no effect.''
\end{example}

\begin{applicationbox}
\textbf{A/B testing: hypothesis testing at internet scale.} Every major website runs
hypothesis tests constantly, under the name \emph{A/B testing}. To decide whether a new
checkout button (variant B) beats the current one (variant A), you randomly assign
incoming users to A or B and measure a conversion rate for each. This is exactly a
two-sample test of proportions:
\[
  H_0\!:\ p_A = p_B \quad\text{vs}\quad H_1\!:\ p_A \ne p_B,\qquad
  Z = \frac{\hat p_B - \hat p_A}{\sqrt{\hat p(1-\hat p)\left(\tfrac{1}{n_A}+\tfrac{1}{n_B}\right)}},
\]
where $\hat p$ pools both groups under the null. \emph{Randomization} is what makes it
work: by assigning users at random, A/B testing turns observational guesswork into
a genuine controlled experiment, so a significant difference can be \emph{causally}
attributed to the button, not to confounders. Companies run thousands of these
simultaneously --- which is exactly why the base-rate and multiple-comparisons traps
below bite hard, and why disciplined teams fix $\alpha$, compute required sample size
\emph{in advance} (power analysis, Example~\ref{ch32:ex-power}), and resist the urge to
``peek.''
\end{applicationbox}

\begin{trapbox}
\textbf{Peeking and the optional-stopping trap.} If you watch an A/B test's p-value live
and stop the moment it dips below $0.05$, your true Type I error is far above $5\%$ ---
sometimes $30\%$ or more. The reason: the p-value wanders randomly over time, and a
wandering quantity will eventually cross any fixed line by chance even if $H_0$ is true.
You have effectively run many tests and reported only the lucky one. The fix is to fix
the sample size in advance, or to use sequential methods (alpha-spending, sequential
probability ratio tests) explicitly designed to permit valid peeking.
\end{trapbox}

\begin{applicationbox}
\textbf{Is model B really better than model A?} When a machine-learning team reports
that a new model beats the old one by $0.4$ accuracy points on a benchmark, the very
first question is the one this chapter is built around: \emph{is that gap real, or test-set
noise?} A held-out test set is itself a finite sample, so accuracy is an estimate with a
standard error. For a test set of $m$ examples and accuracy $\hat a$, the standard error
is $\sqrt{\hat a(1-\hat a)/m}$, and a paired test (the \emph{McNemar test}, a chi-square
on the examples where the two models \emph{disagree}) checks whether the difference
exceeds chance. With $m=10{,}000$ the standard error of an accuracy near $0.9$ is about
$\sqrt{0.9\cdot 0.1/10000}=0.003$, so a $0.4$-point gain ($0.004$) is barely one standard
error --- not significant. Leaderboards that crown winners by hundredths of a percent are
usually reporting noise. The same $\sqrt n$ humility that tames pollsters tames
benchmark hype.
\end{applicationbox}

\begin{trapbox}
\textbf{Multiple comparisons.} Test $20$ independent true nulls at $\alpha=0.05$ and the
chance that \emph{at least one} comes up ``significant'' is $1-(0.95)^{20}=0.64$. Test
enough things and you will find ``significant'' nonsense. The cartoon of the scientist
who tests $20$ jelly-bean colors for acne and triumphantly reports the one with $p<0.05$
is not a joke about bad scientists --- it is a mathematical inevitability. Corrections
like Bonferroni (use $\alpha/m$ for $m$ tests) or controlling the false discovery rate
keep the error rate where you want it.
\end{trapbox}

\section{Putting it together: a complete worked study}
\label{ch32:worked}

\begin{walkbox}
A clinic claims its weight-loss program produces an average loss of more than $5$ kg
over $12$ weeks. You enroll $n=16$ volunteers and record their losses; the sample mean
is $\bar X = 6.2$ kg with sample standard deviation $s = 3.0$ kg. Does the data support
the claim, and how big might the true effect be?
\end{walkbox}

\textbf{Step 1 --- frame the hypotheses.} The claim is ``more than $5$,'' so this is
one-sided: $H_0\!:\mu \le 5$ versus $H_1\!:\mu > 5$. We test at the boundary $\mu_0=5$.

\textbf{Step 2 --- choose the test.} $\sigma$ unknown, small sample, testing a mean:
one-sample, one-sided \textbf{t-test} with $n-1 = 15$ degrees of freedom.

\textbf{Step 3 --- compute the statistic.}
\[
  \mathrm{SE} = \frac{s}{\sqrt n} = \frac{3.0}{4} = 0.75,\qquad
  T = \frac{\bar X - \mu_0}{\mathrm{SE}} = \frac{6.2 - 5}{0.75} = 1.60.
\]

\textbf{Step 4 --- the p-value.} For $t_{15}$, $\P(T \ge 1.60) \approx 0.065$ (one-sided).

\textbf{Step 5 --- decide and report.} At $\alpha=0.05$, since $0.065 > 0.05$,
we \emph{fail to reject} $H_0$: the data do not quite establish that the true mean loss
exceeds $5$ kg, though they lean that way. The $95\%$ confidence interval for $\mu$ is
\[
  6.2 \pm t_{15,0.025}\,\mathrm{SE} = 6.2 \pm 2.131 \times 0.75 = 6.2 \pm 1.60 = (4.6,\ 7.8).
\]
The interval includes $5$ --- consistent with failing to reject --- but also stretches
to $7.8$, so a clinically meaningful effect remains very plausible. The right conclusion
is not ``the program does not work'' but ``promising; the study was too small to confirm
the specific $5$-kg claim --- run a larger trial.'' This is the disciplined, careful
voice the whole chapter has been building toward.

\section{Summary}
\label{ch32:summary}

\begin{aestheticbox}
Every tool in this chapter is a reading of one instrument: the sampling distribution,
the predicted wobble of an estimate across hypothetical repetitions, made Normal by the
CLT and scaled by the standard error $\sigma/\sqrt n$.
\begin{itemize}
  \item A \textbf{confidence interval} inverts the wobble into a range; its ``$95\%$'' is
  long-run \emph{coverage of the procedure}, not a probability about your one interval.
  \item A \textbf{hypothesis test} asks whether the data are surprising under a null; the
  \textbf{p-value} measures that surprise --- and is \emph{not} $\P(H_0\given$ data$)$,
  not effect size, and not a sacred line at $0.05$.
  \item Use \textbf{z} when $\sigma$ is known, \textbf{t} when it is estimated (wider,
  heavier tails, the price of not knowing the spread).
  \item \textbf{Type I} ($\alpha$, false alarm) and \textbf{Type II} ($\beta$, missed
  effect) trade off at fixed $n$; \textbf{power} $=1-\beta$ rises only with more data.
  \item Intervals and tests are \textbf{dual}; A/B testing is hypothesis testing with
  randomization; peeking, multiple comparisons, and the base-rate trap are where sound
  practice is won or lost.
\end{itemize}
The throughline is humility wearing the clothes of mathematics: never a number alone,
always a number with its uncertainty attached.
\end{aestheticbox}

\section{Exercises}
\label{ch32:exercises}

\begin{exercise}[Margin and sample size]
A survey of $n=400$ people estimates a proportion $\hat p = 0.35$. (a) Give the $95\%$
confidence interval. (b) How large must $n$ be to shrink the margin of error to $\pm
0.0234$? (c) Explain in one sentence why the answer to (b) is roughly four times (a)'s $n$
when you halve the margin.
\end{exercise}

\begin{exercise}[z versus t]
You have $n=9$ observations with $\bar X = 20$ and $s = 6$. (a) Build the $95\%$
$t$-interval (use $t_{8,0.025}=2.306$). (b) Build the (incorrect) interval you would get
by using $1.96$. (c) By what percentage does using $z$ understate the half-width, and
why does that misrepresent your precision?
\end{exercise}

\begin{exercise}[Coin test]
A coin gives $44$ heads in $80$ flips. (a) Compute the two-sided z-test statistic and
p-value for $H_0\!: p=0.5$. (b) State your conclusion at $\alpha=0.05$. (c) Give the
$95\%$ confidence interval for $p$ and confirm it is consistent with your test via
duality.
\end{exercise}

\begin{exercise}[Interpreting the interval --- conceptual]
A study reports a $95\%$ confidence interval of $(2.1, 4.7)$ for a treatment effect. For
each statement, say whether it is correct and why: (a) ``There is a $95\%$ probability
the true effect is between $2.1$ and $4.7$.'' (b) ``If we repeated the study many times,
about $95\%$ of the intervals would contain the true effect.'' (c) ``We can reject
$H_0\!:\text{effect}=0$ at the $5\%$ level.'' (d) ``$95\%$ of treated subjects had an
effect between $2.1$ and $4.7$.''
\end{exercise}

\begin{exercise}[The p-value is not what you think]
A test on a true null ($H_0$ holds) produces $p=0.03$, and a colleague concludes ``there
is only a $3\%$ chance the null is true.'' (a) Identify the precise error. (b) Explain
what the $3\%$ actually refers to. (c) Using the base-rate paradox, sketch why even a
genuine $p=0.03$ can correspond to a much higher probability that $H_0$ is true when
true effects are rare.
\end{exercise}

\begin{exercise}[Power]
A one-sided z-test at $\alpha=0.05$ tests $H_0\!:\mu=0$ against $H_1\!:\mu=\delta>0$ with
$\sigma=10$ and $n=25$. (a) Find the rejection cutoff for $\bar X$. (b) Compute the power
when $\delta=4$. (c) What $n$ would raise the power to $0.9$ for $\delta=4$? (Use
$\Phi(1.28)=0.90$.)
\end{exercise}

\begin{exercise}[Two-sample A/B test]
Variant A converts $180$ of $3000$ visitors; variant B converts $230$ of $3000$. (a)
Compute $\hat p_A$, $\hat p_B$, and the pooled $\hat p$. (b) Compute the two-sided
z-statistic and p-value for $H_0\!: p_A=p_B$. (c) State the business conclusion, and
name one trap (peeking, multiple comparisons, base rate) that could still make this
result misleading.
\end{exercise}

\begin{exercise}[Paired versus unpaired --- conceptual]
Twelve runners are timed before and after a training regimen. Explain why a
\emph{paired} t-test on the within-runner differences is far more powerful than an
unpaired two-sample test treating ``before'' and ``after'' as independent groups. What
source of variation does pairing remove?
\end{exercise}

\begin{exercise}[Duality, proved]
For the known-$\sigma$ case, prove from the definitions that a value $\mu_0$ lies in the
$(1-\alpha)$ confidence interval $\bar X_n \pm z_{\alpha/2}\,\sigma/\sqrt n$ if and only
if the two-sided level-$\alpha$ z-test of $H_0\!:\mu=\mu_0$ fails to reject. (This is the
interval--test duality of Section~\ref{ch32:duality}.)
\end{exercise}

\begin{exercise}[Why $n-1$? --- proof]
Let $X_1,\dots,X_n$ be i.i.d.\ with variance $\sigma^2$. Show that
$\E\!\left[\sum_{i=1}^n (X_i - \bar X_n)^2\right] = (n-1)\sigma^2$, and hence that
dividing by $n-1$ (not $n$) makes $s^2$ unbiased. \emph{Hint:} write $X_i - \bar X_n =
(X_i - \mu) - (\bar X_n - \mu)$, expand the square, and use $\Var(\bar X_n)=\sigma^2/n$.
\end{exercise}

\begin{exercise}[Significance versus importance]
A study of $n=1{,}000{,}000$ users finds that a redesign increases average session length
from $300.0$ to $300.5$ seconds, with $s=120$ seconds, $p < 10^{-4}$. (a) Verify the
result is ``highly significant.'' (b) Compute the $95\%$ confidence interval for the
increase. (c) Argue, using the interval, why ``statistically significant'' here does not
mean ``worth shipping.''
\end{exercise}
